%  =============================================================================
%  Datei-Hinweise & Konfiguration: siehe config.md
%  -> Abschnitt "anhang/ki-prompts/tool-gemini.tex"
%  =============================================================================

\section{Google Gemini (Flash~3.5)}
\label{sec:prompts-gemini}

\kiwerkzeug{Google Gemini, Modell Flash 3.5}{Stand: Mai 2026}{https://gemini.google.com}
% Eintrag im KI-Verzeichnis (Leitfaden Kap. 6.2.4) erzwingen mittels \nocite.
\nocite{geminiFlash2026}

\subsection{Literaturrecherche}
\label{sec:prompts-gemini-literatur}
\begin{enumerate}
    \item \enquote{Nenne einschlägige, zitierfähige Quellen (Monografien,
            Normen, begutachtete Fachartikel) zur BPMN-2.0-Prozessmodellierung
            und gib jeweils vollständige bibliografische Angaben an: [\dots]}
            \par\hfill\textit{Kontext: \Cref{cha:grundlagen}}

    \item \enquote{Fasse den aktuellen Forschungsstand zur
            Geschäftsprozessmodellierung im ERP-Umfeld zusammen und benenne
            die zentralen Referenzwerke: [\dots]}
            \par\hfill\textit{Kontext: \Cref{cha:grundlagen}}
\end{enumerate}

\subsection{Erstellung des Abkürzungsverzeichnisses}
\label{sec:prompts-gemini-abkuerzungen}
\begin{enumerate}
    \item \enquote{Extrahiere alle Abkürzungen aus dem folgenden Text und
            erstelle ein alphabetisch sortiertes Abkürzungsverzeichnis mit
            jeweiliger Langform: [\dots]}
            \par\hfill\textit{Kontext: gesamte Arbeit}
    \item \enquote{Überprüfe das vorhandene Abkürzungsverzeichnis auf
            Vollständigkeit und Konsistenz und ergänze fehlende Einträge: [\dots]}
            \par\hfill\textit{Kontext: gesamte Arbeit}
\end{enumerate}

%  =============================================================================
%  Ende
%  =============================================================================
